\chapter{FUNKTOR $\mathbf{\emph{K}_0}$}

Sekarang kita meninjau secara singkat apa yang disebut teori-$K$ untuk aljabar-$C^*$. Barangkali pengantar terbaik untuk pokok bahasan ini adalah~\cite{RordamLL:2000},
yang beberapa bagian elementernya diikuti dengan cukup saksama dalam catatan ini. Pengantar lain yang sangat mudah dibaca adalah~\cite{Wegge-Olsen:1993}.

Diberikan suatu aljabar-$C^*$ $A$, kita terutama akan memperhatikan dua grup yang dikenal sebagai $K_0(A)$ dan~$K_1(A)$. Bab ini membahas
grup yang pertama, yakni grup proyeksi di $A$ dengan operasi penjumlahan. Tentu saja, jika kita ingin menjumlahkan proyeksi sembarang, kita langsung
menghadapi kesulitan serius. Sebelumnya telah kita tunjukkan (dalam proposisi~\ref{prop_sum_projs}) bahwa agar proyeksi dapat dijumlahkan, proyeksi-proyeksi itu harus komutatif
(dan hasil kalinya nol)! Penyelesaian dilema ini mencerminkan cara berpikir yang lazim dalam teori-$K$ secara umum. Jika dua proyeksi tidak
komutatif, singkirkan penghalang yang mencegah keduanya menjadi komutatif. Jika tidak cukup ruang bagi keduanya untuk saling melewati, berilah keduanya lebih banyak
ruang. Jangan bersikeras memandang keduanya sebagai makhluk yang mencoba hidup dalam dunia matriks $1 \times 1$ berentri di~$A$ yang sesak tanpa harapan.
Biarkan keduanya, misalnya, menjelajahi dunia matriks $2 \times 2$ yang jauh lebih lapang. % SOURCE-CORRECTION: CH17-C001
Untuk mewujudkan gagasan ini secara teknis, kita mencari relasi ekuivalensi $\sim$ yang mengidentifikasi proyeksi $p$ di $A$ dengan matriks $\begin{bmatrix} p & \vc 0 \\ \vc 0 & \vc 0 \end{bmatrix}$
dan $\begin{bmatrix} \vc 0 & \vc 0 \\ \vc 0 & p \end{bmatrix}$. Sebagai motivasi heuristik saja---bukan sebagai konsekuensi suatu kongruensi perkalian yang telah dibuktikan---hambatan itu tampak hilang: jika $p$ dan $q$ proyeksi di~$A$, maka % SOURCE-CORRECTION: CH17-C002
    \[ pq \sim \begin{bmatrix} p & \vc 0 \\ \vc 0 & \vc 0 \end{bmatrix} \begin{bmatrix} \vc 0 & \vc 0 \\ \vc 0 & q \end{bmatrix}  \\
                   = \begin{bmatrix} \vc 0 & \vc 0 \\ \vc 0 & \vc 0 \end{bmatrix}           \\
                   =  \begin{bmatrix} \vc 0 & \vc 0 \\ \vc 0 &  q \end{bmatrix}\begin{bmatrix} p & \vc 0 \\ \vc 0 & \vc 0 \end{bmatrix}  \\
                   \sim qp\,. \]
Dalam gambaran heuristik ini, $p$ dan $q$ komutatif modulo relasi ekuivalensi $\sim$, dan kita dapat mendefinisikan penjumlahan dengan sesuatu seperti
$p \oplus q = \begin{bmatrix} p & \vc 0 \\ \vc 0 & q \end{bmatrix}$.

Cukup jelas bahwa perangkat sederhana di atas tidak akan langsung menghasilkan sesuatu yang menyerupai grup Abelian, tetapi ini merupakan suatu permulaan. Jika Anda menduga
bahwa pada akhirnya konstruksi $K_0(A)$ mungkin agak rumit, dugaan Anda sepenuhnya benar.




\section{Relasi Ekuivalensi pada Proyeksi}

Kita mendefinisikan beberapa relasi ekuivalensi, yang semuanya sesuai untuk proyeksi dalam aljabar-$C^*$ beridentitas.

\begin{defn} Dalam suatu aljabar beridentitas, elemen $a$ dan $b$ disebut
 \index{serupa!elemen suatu aljabar}%
\df{serupa} jika terdapat elemen invertibel $s$ sedemikian sehingga $b = sas^{-1}$.
 \index{<binrelequivs@$a \sim_s b$ (keserupaan elemen suatu aljabar)}%
Dalam hal ini kita menulis $a \sim_s b$.
\end{defn}

\begin{prop} Relasi keserupaan $\sim_s$ yang didefinisikan di atas merupakan relasi ekuivalensi pada elemen-elemen suatu aljabar beridentitas.
\end{prop}

\begin{defn} Dalam suatu aljabar-$C^*$, elemen $a$ dan $b$ disebut
 \index{uniter!ekuivalensi}%
\df{ekuivalen secara uniter} jika terdapat suatu elemen uniter $u \in \wt A$ sedemikian sehingga $b = uau^*$. % SOURCE-CORRECTION: CH17-C003
 \index{<binrelequivu@$a \sim_u b$ (ekuivalensi uniter)}%
Dalam hal ini kita menulis $a \sim_u b$.
\end{defn}

\begin{prop} Relasi ekuivalensi uniter $\sim_u$ yang didefinisikan di atas merupakan relasi ekuivalensi pada elemen-elemen suatu aljabar-$C^*$.
\end{prop}

\begin{prop}\label{0060224fa} Elemen $a$ dan $b$ dari suatu aljabar-$C^*$ beridentitas $A$ ekuivalen secara uniter jika dan hanya jika
terdapat elemen $u \in \ofml U(A)$ sedemikian sehingga $b = uau^*$.
\end{prop}

\begin{proof}[\emph{Petunjuk untuk bukti}] Andaikan terdapat $v \in \ofml U(\wt A)$ sedemikian sehingga $b = vav^*$. Berdasarkan proposisi~\ref{0015213},
terdapat $u \in A$ dan $\alpha \in \C$ sedemikian sehingga $v = u + \alpha \vc j$ (dengan $\vc j = \vc 1_{\wt A} - \vc 1_A$). Tunjukkan
bahwa $\abs\alpha = 1$, $u$ uniter, dan $b = uau^*$.

Untuk arah sebaliknya, andaikan terdapat $u \in \ofml U(A)$ sedemikian sehingga $b = uau^*$. Misalkan $v = u + \vc j$. \ns
\end{proof}

\begin{prop}[Dekomposisi polar]\label{0060141} Untuk setiap elemen invertibel $s$ dari suatu
 \index{dekomposisi polar!elemen invertibel aljabar-$C^*$}%
 \index{dekomposisi!polar}%
aljabar-$C^*$ beridentitas, terdapat suatu elemen uniter $\omega(s)$ dalam aljabar tersebut sedemikian sehingga % SOURCE-CORRECTION: CH17-C004
   \[   s = \omega(s) \abs s\,. \]
\end{prop}

\begin{proof}[\emph{Petunjuk untuk bukti}] Gunakan Akibat~\ref{00183313}.   \ns
\end{proof}

\begin{prop}\label{K002021} Jika $A$ suatu aljabar-$C^*$ beridentitas, maka fungsi $\omega\colon \inv A \sto \ofml U(A)$
yang didefinisikan dalam proposisi sebelumnya kontinu.
\end{prop}

\begin{proof}[\emph{Petunjuk untuk bukti}] Gunakan proposisi~\ref{000646fa} dan~\ref{001449}.   \ns
\end{proof}

\begin{prop}\label{K002024} Misalkan $s = u\abs s$ dekomposisi polar dari elemen invertibel $s$ dalam aljabar-$C^*$ beridentitas~$A$.
Maka $u \sim_h s$ dalam $\inv A$.
\end{prop}

\begin{proof}[\emph{Petunjuk untuk bukti}] Untuk $0 \le t \le 1$, misalkan $c_t = u(t\abs s + (1 - t)\vc 1)$. Simpulkan dari
proposisi~\ref{0018301} bahwa terdapat $\epsilon \in (0,1]$ sedemikian sehingga $\abs s \ge \epsilon \vc 1$. Gunakan proposisi yang sama
untuk menunjukkan bahwa $t\abs s + (1 - t)\vc 1$ invertibel untuk setiap $t \in [0,1]$.   \ns
\end{proof}

\begin{prop}\label{K002027} Misalkan $u$ dan $v$ elemen-elemen uniter dalam aljabar-$C^*$ beridentitas~$A$. Jika $u \sim_h v$ dalam $\inv A$,
maka $u \sim_h v$ dalam~$\ofml U(A)$.
\end{prop}

\begin{prop}\label{K002031} Misalkan $u_1$, $u_2$, $u_3$, dan $u_4$ elemen-elemen uniter dalam aljabar-$C^*$ beridentitas~$A$.
Jika $u_1 \sim_h u_2$ dalam $\ofml U(A)$ dan $u_3 \sim_h u_4$ dalam $\ofml U(A)$, maka $u_1u_3 \sim_h u_2u_4$ dalam~$\ofml U(A)$.
\end{prop}

\begin{exam}\label{0060116} Misalkan $h$ suatu elemen swaadjoin dari aljabar-$C^*$ beridentitas~$A$. Maka $\exp(ih)$ uniter dan
homotop dengan~$\vc 1$ dalam~$\ofml U(A)$.
\end{exam}

\begin{proof}[\emph{Petunjuk untuk bukti}] Pertimbangkan lintasan $c\colon [0,1] \sto \ofml U(A)\colon t \mapsto \exp(ith)$. Gunakan % SOURCE-CORRECTION: CH17-C005
proposisi~\ref{001449}.  \ns
\end{proof}

\begin{notn} Jika $p$ dan $q$ proyeksi dalam suatu aljabar-$C^*$ $A$, kita menulis $p \sim q$
 \index{ekuivalensi Murray--von Neumann}%
 \index{ekuivalensi!Murray--von Neumann}%
 \index{<binrelequiv@$p \sim q$ (ekuivalensi Murray--von Neumann)}%
($p$ \df{ekuivalen Murray--von Neumann} dengan~$q$) jika terdapat elemen $v \in A$
sedemikian sehingga $v^*v = p$ dan $vv^* = q$. Perhatikan bahwa $v$ demikian secara otomatis merupakan isometri
parsial. Kita akan menyebutnya
 \index{realisasi ekuivalensi Murray--von Neumann}%
 \index{ekuivalensi!Murray--von Neumann!realisasi}%
 \index{Murray--von Neumann!ekuivalensi!realisasi}%
\emph{isometri parsial yang merealisasikan ekuivalensi}.
\end{notn}

\begin{prop}\label{0060216} Relasi ekuivalensi Murray--von Neumann $\sim$ merupakan
relasi ekuivalensi pada keluarga $\fml P(A)$ dari proyeksi-proyeksi dalam suatu aljabar-$C^*$~$A$.
\end{prop}

\begin{prop}\label{0060144} Misalkan $a$ dan $b$ elemen-elemen swaadjoin dari suatu aljabar-$C^*$ beridentitas. Jika $a \sim_s b$, maka $a \sim_u b$.
Bahkan, jika $b = sas^{-1}$ dan $s = u\abs s$ adalah dekomposisi polar dari $s$, maka $b = uau^*$.
\end{prop}

\begin{proof}[\emph{Petunjuk untuk bukti}] Andaikan terdapat elemen invertibel $s$ sedemikian sehingga $b = sas^{-1}$. Misalkan $s = u\abs s$
dekomposisi polar dari~$s$. Tunjukkan bahwa $a$ komutatif dengan ${\abs s}^2$, sehingga juga komutatif dengan setiap elemen dalam aljabar $C^*(\vc 1,{\abs s}^2)$.
Secara khusus, $a$ komutatif dengan~${\abs s}^{-1}$. Dari sini diperoleh $uau^* = b$.   \ns
\end{proof}

\begin{prop}\label{0060221} Jika $p$ dan $q$ proyeksi dalam suatu aljabar-$C^*$, maka
   \[ p \sim_h q \implies p \sim_u q \implies p \sim q\,. \]
\end{prop}

\begin{proof}[\emph{Petunjuk untuk bukti}] Dalam petunjuk ini $\vc 1 = \vc 1_{\wt A}$. Untuk implikasi pertama, tunjukkan bahwa tanpa mengurangi keumuman
kita dapat mengandaikan $\norm{p - q} < \frac12$. Misalkan $s = pq + (\vc 1 - p)(\vc 1 - q)$. Buktikan bahwa $p \sim_s q$. Untuk
itu, tuliskan $s - \vc 1$ sebagai $p(q - p) + (\vc 1 - p)((\vc 1 - q) - (\vc 1 - p))$ dan gunakan akibat~\ref{cor_Neumann_series}.
Kemudian gunakan proposisi~\ref{0060144}.

Untuk membuktikan implikasi kedua, perhatikan bahwa jika $upu^* = q$ untuk suatu elemen uniter dalam~$\wt A$, maka $up$ merupakan isometri parsial
dalam~$A$.  \ns
\end{proof}

Secara umum, kebalikan dari implikasi kedua, $p \sim q \implies p \sim_u q$, dalam proposisi sebelumnya tidak berlaku (lihat % SOURCE-CORRECTION: CH17-C024
contoh~\ref{0060233} di bawah). Namun, untuk proyeksi $p$ dan $q$ dalam suatu aljabar-$C^*$ beridentitas, jika berlaku baik $p \sim q$ maupun
$\vc 1 - p \sim \vc 1 - q$, maka kita dapat menyimpulkan $p \sim_u q$.

\begin{prop}\label{0060224} Misalkan $p$ dan $q$ proyeksi dalam suatu aljabar-$C^*$ beridentitas~$A$. Maka $p \sim_u q$ jika dan hanya jika
$p \sim q$ dan $\vc 1 - p \sim \vc 1 - q$.
\end{prop}

\begin{proof}[\emph{Petunjuk untuk bukti}] Untuk arah sebaliknya, andaikan isometri parsial $v$ merealisasikan ekuivalensi $p \sim q$
dan $w$ merealisasikan $\vc 1 - p \sim \vc 1 - q$. Pertimbangkan elemen $u = v + w + \vc j$ dalam~$\wt A$.   \ns
\end{proof}

\begin{defn} Suatu elemen $s$ dari aljabar-$C^*$ beridentitas disebut
 \index{isometri!dalam aljabar-$C^*$ beridentitas}%
\df{isometri} jika $s^*s = \vc 1$.
\end{defn}

\begin{exer} Jelaskan mengapa terminologi dalam definisi sebelumnya masuk akal.
\end{exer}

Tidak sulit untuk melihat bahwa kebalikan dari implikasi kedua dalam proposisi~\ref{0060221} pada umumnya gagal.

\begin{exam}\label{0060233} Jika $p$ dan $q$ proyeksi dalam suatu aljabar-$C^*$, maka
   \[ p \sim q \,\,\,\,\not\!\!\!\!\implies p \sim_u q\,. \]
Sebagai contoh, jika $s$ merupakan isometri tak uniter (seperti geser unilateral), maka $s^*s \sim ss^*$, tetapi $s^*s \nsim_u ss^*$.
\end{exam}

\begin{proof}[\emph{Petunjuk untuk bukti}] Buktikan, dan ingatlah, bahwa tidak ada proyeksi tak nol yang dapat ekuivalen Murray--von Neumann
dengan proyeksi nol.   \ns
\end{proof}

\begin{rem} Kebalikan dari implikasi pertama, $p \sim_u q \implies p \sim_h q$, dalam proposisi~\ref{0060221} juga tidak berlaku % SOURCE-CORRECTION: CH17-C025
secara umum untuk proyeksi dalam suatu aljabar-$C^*$. Namun, contoh yang menggambarkan gejala ini tidak mudah ditemukan. Untuk melihat
apa yang terlibat, lihat~\cite{RordamLL:2000}, contoh 2.2.9 dan 11.3.4.
\end{rem}

Alangkah baiknya jika implikasi-implikasi dalam proposisi~\ref{0060221} dapat dibalik. Namun, seperti telah kita lihat, hal itu tidak terjadi.
Salah satu cara menghadapi fakta yang membandel, sebagaimana dicatat pada awal bab ini, adalah memberi objek-objek matematika
yang sedang kita kaji ruang yang lebih luas untuk bergerak. Beralihlah ke matriks. Proposisi~\ref{0060236} dan~\ref{0060244} merupakan contoh
cara kerja teknik ini. Dalam arti tertentu, kita dapat membuat implikasi-implikasi dalam~\ref{0060221} berbalik arah.

\begin{notn} Jika $a_1$, $a_2$, \dots $a_n$ elemen-elemen dari suatu aljabar-$C^*$ $A$, maka
 \index{diagonal@$\diag(a_1,\dots,a_n)$ (matriks diagonal dalam $\M nA$)}%
$\diag(a_1,a_2,\dots,a_n)$ adalah matriks diagonal dalam $\M nA$ yang diagonal utamanya terdiri atas
elemen-elemen $a_1$, \dots, $a_n$. Notasi ini juga kita gunakan untuk matriks blok. Sebagai
contoh, jika $a$ matriks $m \times m$ dan $b$ matriks $n \times n$, maka
$\diag(a,b)$ adalah matriks $(m+n) \times (m+n)$ $\begin{bmatrix} a & \vc 0 \\ \vc 0 & b
\end{bmatrix}$. (Tentu saja, $\vc 0$ di sudut kanan atas matriks ini adalah matriks $m \times n$ yang semua
entrinya nol, sedangkan $\vc 0$ di sudut kiri bawah adalah matriks $n \times m$ yang setiap entrinya nol.)
\end{notn}

\begin{prop}\label{0060236} Misalkan $p$ dan $q$ proyeksi dalam suatu aljabar-$C^*$~$A$. Maka
   \[ p \sim q \implies \diag(p,\vc 0) \sim_u \diag(q,\vc 0) \text{ dalam $\M 2A$.} \]
\end{prop}

\begin{proof}[\emph{Petunjuk untuk bukti}] Misalkan $v$ isometri parsial dalam $A$ yang merealisasikan ekuivalensi Murray--von Neumann $p \sim q$.
Pertimbangkan matriks $u = \begin{bmatrix} v  &  \vc 1_{\wt A} - q  \\  \vc 1_{\wt A} - p  &  v^* \end{bmatrix}$.  \ns
\end{proof}

Ingat dari akibat~\ref{001311005fa} bahwa spektrum setiap elemen uniter dari suatu aljabar-$C^*$ beridentitas terletak pada lingkaran satuan.
Hebatnya, jika spektrumnya bukan seluruh lingkaran, elemen itu homotop dengan~$\vc 1$ dalam~$\ofml U(A)$.

\begin{prop}\label{0060121} Misalkan $u$ suatu elemen uniter dalam aljabar-$C^*$ beridentitas. Jika $\sigma(u) \ne \T$, maka $u \sim_h \vc 1$ dalam~$\ofml U(A)$.
\end{prop}

\begin{proof}[\emph{Petunjuk untuk bukti}] Pilih $\theta \in \R$ sedemikian sehingga $\exp(i\theta) \notin \sigma(u)$. Maka terdapat fungsi kontinu yang tunggal
   \[ \phi\colon \sigma(u) \sto (\theta, \theta + 2\pi)\colon \exp(it) \mapsto t\,. \]
Misalkan $h = \phi(u)$ dan gunakan contoh~\ref{0060116}.   \ns
\end{proof}

\begin{exam}\label{0060129a} Jika $A$ suatu aljabar-$C^*$ beridentitas, maka $\begin{bmatrix} \vc 0  &  \vc 1  \\  \vc 1  &  \vc 0  \end{bmatrix} \sim_h
\begin{bmatrix} \vc 1  &  \vc 0  \\ \vc 0  & \vc 1 \end{bmatrix}$ dalam~$\ofml U(\M 2A)$.
\end{exam}

Matriks $J = \begin{bmatrix} \vc 0  &  \vc 1  \\ \vc 1 & \vc 0 \end{bmatrix}$ sangat berguna jika dipakai bersama
proposisi~\ref{K002031} untuk membentuk homotopi antara matriks-matriks dalam $\ofml U(A)$. Dalam memeriksa beberapa contoh
berikutnya, matriks ini sangat membantu. Perhatikan bahwa mengalikan matriks $2 \times 2$ dari kanan dengan $J$ menukar kolom-kolomnya; mengalikan
dari kiri dengan $J$ menukar baris-barisnya; dan mengalikan dari kiri serta kanan dengan $J$ menukar elemen-elemen pada kedua diagonal.

\begin{exam}\label{0060129b} Jika $u$ dan $v$ elemen-elemen uniter dalam suatu aljabar-$C^*$ beridentitas~$A$, maka
   \[ \diag(u,v) \sim_h \diag(v,u) \]
dalam~$\ofml U(\M 2A)$.
\end{exam}

\begin{exam}\label{0060129c} Jika $u$ dan $v$ elemen-elemen uniter dalam suatu aljabar-$C^*$ beridentitas $A$, maka
  \[ \diag(u,v) \sim_h \diag(uv, \vc 1) \]
dalam $\ofml U(\M 2A)$.
\end{exam}

\begin{exam}\label{0060129d} Jika $u$ dan $v$ elemen-elemen uniter dalam suatu aljabar-$C^*$ beridentitas~$A$, maka
   \[ \diag(uv,\vc 1) \sim_h \diag(vu, \vc 1) \]
dalam~$\ofml U(\M 2A)$.
\end{exam}

\begin{prop}\label{0060244} Misalkan $p$ dan $q$ proyeksi dalam suatu aljabar-$C^*$~$A$. Maka
   \[ p \sim_u q \implies \diag(p,\vc 0) \sim_h \diag(q,\vc 0) \text{ dalam $\ofml P(\M 2A)$.} \]
\end{prop}

\begin{exam}\label{0060247fa} Dua proyeksi dalam $\mathbf M_n$ ekuivalen Murray--von Neumann jika dan hanya jika keduanya memiliki teras yang sama,
yang pada gilirannya ekuivalen dengan jangkauan keduanya berdimensi sama.
\end{exam}

\begin{proof}[\emph{Petunjuk untuk bukti}] Gunakan proposisi~\ref{tr0121}--\ref{tr0124} dan~\ref{pi0117}--\ref{pi0121}.   \ns
\end{proof}

Dalam contoh~\ref{000533a} kita memperkenalkan aljabar matriks $n \times n$ $\M nA$, dengan $A$ suatu aljabar. Jika $A$ suatu
aljabar-$*\,$ dan $n \in \N$, kita dapat memperlengkapi $\M nA$ dengan involusi secara alami. Involusi itu didefinisikan, seperti yang dapat diduga,
sebagai analog dari ``transposisi konjugat'': jika $\vc a \in \M nA$, maka
   \[ \vc a^* = \bigl[a_{ij}\bigr]^* := \bigl[{a_{ji}}^*\bigr]. \]

Jika $A$ suatu aljabar-$C^*$, kita dapat memperkenalkan norma pada $\M nA$ yang membuatnya juga menjadi aljabar-$C^*$.
 \index{Ma@$\M nA$!sebagai aljabar-$C^*$}%
 \index{C*@aljabar-$C^*$!$\M nA$ sebagai suatu}%
Pilih representasi setia $\phi\colon A \sto \ofml B(H)$ dari $A$, dengan $H$ suatu ruang Hilbert. Untuk setiap $n \in \N$
definisikan
   \[ \phi_n\colon \M nA \sto \ofml B(H^n)\colon \bigl[a_{ij}\bigr] \mapsto \bigl[ \phi(a_{ij})\bigr] \]
dengan $H^n$ jumlah langsung $n$ rangkap dari $H$ dengan dirinya sendiri. Kemudian definisikan norma suatu elemen $\vc a \in \M nA$ sebagai
norma operator $\phi_n(\vc a) \in \ofml B(H^n)$. Yakni, $\norm{\vc a} := \norm{\phi_n(\vc a)}$.

Norma ini merupakan norma-$C^*$ pada $\M nA$. Norma tersebut tidak bergantung pada representasi khusus yang kita pilih berkat
ketunggalan norma-$C^*$ (lihat akibat~\ref{00152144}).

\begin{prop} Misalkan $A$ suatu aljabar-$C^*$ dan $\vc a \in \M nA$. Maka
   \[ \max\{\norm{a_{ij}}\colon 1 \le i,j \le n\} \le \norm{\vc a} \le \sum_{i,j = 1}^n \norm{a_{ij}}. \]
\end{prop}

\begin{proof}[\emph{Petunjuk untuk bukti}] Karena setiap representasi setia merupakan isometri (lihat proposisi~\ref{00152181}),
Anda dapat menyederhanakan persoalan dengan memandang representasi $\phi\colon A \sto \ofml B(H)$ yang dibahas di atas sebagai
pemetaan inklusi.

Untuk ketaksamaan pertama, bagi $v \in H$ definisikan vektor $\vc E_j(v) \in H^n$ sebagai $n$-tupel
yang semua entrinya nol kecuali entri ke-$j$, yang bernilai~$v$. Kemudian periksa bahwa
   \[ \norm{a_{ij}v} \le \norm{\vc a\vc E_j(v)} \le \norm{\vc a} \]
setiap kali $\vc a \in \M nA$, $v \in H$, dan $\norm v \le 1$.

Untuk ketaksamaan kedua, tunjukkan bahwa
   \[ {\norm{\vc a\vc v}}^2 \le \sum_{i=1}^n\biggl(\sum_{j=1}^n \norm{a_{ij}}\biggr)^2 \le \biggl(\sum_{i=1}^n\sum_{j=1}^n \norm{a_{ij}}\biggr)^2 \]
setiap kali $\vc a \in \M nA$, $\vc v \in H^n$, dan $\norm{\vc v} \le 1$.    \ns
\end{proof}


























\section{Semigrup Proyeksi}

\begin{notn} Jika $A$ suatu aljabar-$C^*$ dan $n \in \N$, kita menetapkan
 \index{P@$\fml P_n(A)$, $\fml P_\infty(A)$ (proyeksi dalam aljabar matriks)}%
  \begin{align*}
       \fml P_n(A) &= \fml P(\M nA) \quad\text{ dan} \\
       \fml P_\infty(A) &= \bigcup_{n=1}^\infty \fml P_n(A).
  \end{align*}
\end{notn}

Sekarang kita memperluas ekuivalensi Murray--von Neumann ke matriks-matriks yang ukurannya berbeda.

\begin{notn} Jika $A$ suatu aljabar-$C^*$, misalkan
 \index{Ma@$\M {n,m}A$ (matriks dengan entri dari~$A$)}%
$\M {n,m}A$ menyatakan himpunan matriks $n \times m$ berentri dalam~$A$.
\end{notn}

Selanjutnya kita memperluas ekuivalensi Murray--von Neumann ke matriks proyeksi yang ukurannya tidak harus sama.

\begin{defn}\label{0060316fa} Jika $A$ suatu aljabar-$C^*$, $p \in \fml P_m(A)$, dan $q \in \fml P_n(A)$, kita menetapkan
 \index{ekuivalensi Murray--von Neumann}%
 \index{ekuivalensi!Murray--von Neumann}%
 \index{<binrelequiv@$p \sim q$ (ekuivalensi Murray--von Neumann)}%
$p \sim q$ jika terdapat $v \in \M {n,m}A$ sedemikian sehingga
  \[ v^*v = p \qquad\text{ dan } \qquad  vv^* = q. \]
Kita akan memperluas penggunaan nama \emph{ekuivalensi Murray--von Neumann} untuk relasi baru ini pada~$\fml P_\infty(A)$.
\end{defn}

\begin{prop}\label{0060317} Relasi $\sim$ yang didefinisikan di atas merupakan relasi
ekuivalensi pada~$\fml P_\infty(A)$.
\end{prop}

\begin{defn} Untuk setiap aljabar-$C^*$ $A$, kita mendefinisikan operasi biner $\oplus$ pada $\fml
P_\infty(A)$ dengan
   \[ p \oplus q = \diag(p,q)\,. \]
Jadi, jika $p \in \fml P_m(A)$ dan $q \in \fml P_n(A)$, maka $p \oplus q \in \fml
P_{m+n}(A)$.
\end{defn}

Dalam proposisi berikut, $\vc 0_n$ adalah identitas aditif dalam~$\fml P_n(A)$.

\begin{prop}\label{0060324} Misalkan $A$ suatu aljabar-$C^*$ dan $p \in \fml P_\infty(A)$.
 \index{<sum@$p \oplus q$ (operasi biner dalam $\fml P_\infty(A)$)}%
Maka $p \sim p \oplus \vc 0_n$ untuk setiap $n \in \N$.
\end{prop}

\begin{prop}\label{0060327} Misalkan $A$ suatu aljabar-$C^*$ dan $p$, $p'$, $q$, $q' \in
\fml P_\infty(A)$. Jika $p \sim p'$ dan $q \sim q'$, maka $p \oplus q \sim p' \oplus q'$.
\end{prop}

\begin{prop}\label{0060331} Misalkan $A$ suatu aljabar-$C^*$ dan $p$, $q \in \fml P_\infty(A)$.
Maka $p \oplus q \sim q \oplus p$.
\end{prop}

\begin{prop}\label{0060334} Misalkan $A$ suatu aljabar-$C^*$ dan $p$, $q \in \fml P_n(A)$ untuk
suatu $n$. Jika $p \perp q$, maka $p + q$ suatu proyeksi dalam $\M nA$ dan $p + q \sim p \oplus q$.
\end{prop}

\begin{prop}\label{0060337} Jika $A$ suatu aljabar-$C^*$, maka operasi $\oplus$ pada $\fml P_\infty(A)$ bersifat
asosiatif secara ketat dan komutatif hanya hingga ekuivalensi Murray--von Neumann. % SOURCE-CORRECTION: CH17-C006
\end{prop}

\begin{notn}\label{0060411} Jika $A$ suatu aljabar-$C^*$ dan $p \in \fml P_\infty(A)$, misalkan
 \index{<bracsqr@$\sbsb{[p]}{\fml D}$ (kelas ekuivalensi $\sim$ dari $p$)}%
$\sbsb{[p\,]}{\fml D}$ adalah kelas ekuivalensi yang memuat $p$ dan ditentukan oleh
relasi ekuivalensi~$\sim$. Tetapkan pula
 \index{D@$\fml D(A)$ (himpunan kelas ekuivalensi proyeksi)}%
$\fml D(A) := \{\,\sbsb{[p\,]}{\fml D}\colon p \in \fml P_\infty(A)\,\}$.
\end{notn}

\begin{defn}\label{0060414} Misalkan $A$ suatu aljabar-$C^*$. Definisikan operasi biner $+$ pada
$\fml D(A)$ dengan
 \index{<binop@$\sbsb{[p\,]}{\fml D} + \sbsb{[q\,]}{\fml D}$ (penjumlahan dalam $\fml D(A)$)}%
   \[ \sbsb{[p\,]}{\fml D} + \sbsb{[q\,]}{\fml D} := \sbsb{[p \oplus q\,]}{\fml D} \]
dengan $p$, $q \in \fml P_\infty(A)$.
\end{defn}

\begin{prop}\label{0060417} Operasi $+$ yang didefinisikan dalam~\ref{0060414} terdefinisi dengan baik
dan menjadikan $\fml D(A)$ suatu semigrup komutatif.
\end{prop}

\begin{exam}\label{0060421} Semigrup $\fml D(\C)$ isomorfik dengan semigrup aditif bilangan bulat tak negatif; yaitu, % SOURCE-CORRECTION: CH17-C007
    \[ \fml D(\C) \cong \Z^+ = \{0,1,2,\dots\}\,.  \]
\end{exam}

\begin{proof}[\emph{Petunjuk untuk bukti}] Gunakan contoh~\ref{0060247fa}. \ns
\end{proof}

\begin{exam}\label{0060424} Jika $H$ suatu ruang Hilbert terpisahkan berdimensi tak hingga, maka % SOURCE-CORRECTION: CH17-C008
   \[ \fml D(\ofml B(H)) \cong \Z^+ \cup \{\infty\}. \]
(Gunakan penjumlahan biasa pada $\Z^+$ dan tetapkan $n + \infty = \infty + n = \infty$ setiap kali $n
\in \Z^+ \cup \{\infty\}$.)
\end{exam}

\begin{exam}\label{0060427} Untuk aljabar-$C^*$ $\C \oplus \C$, kita mempunyai
   \[ \fml D(\C \oplus \C) \cong \Z^+ \oplus \Z^+\,. \]
\end{exam}


























\section{Konstruksi Grothendieck}
Dalam konstruksi medan bilangan real, kita dapat menggunakan teknik yang sama untuk beralih dari (semigrup aditif) bilangan asli ke
(grup aditif semua) bilangan bulat seperti yang kita gunakan untuk beralih dari (semigrup multiplikatif) bilangan bulat tak nol ke (grup multiplikatif
semua) bilangan rasional tak nol. Teknik ini disebut \emph{konstruksi Grothendieck}.

\begin{defn}\label{0062111} Misalkan $(S,+)$ suatu semigrup komutatif. Definisikan relasi
 \index{<binrelequiv@$(a,b) \sim (c,d)$ (ekuivalensi Grothendieck)}%
$\sim$ pada $S \times S$ dengan
   \[ (a,b) \sim (c,d) \text{ jika terdapat $k \in S$ sedemikian sehingga } a + d + k = b + c + k. \]
\end{defn}

\begin{prop}\label{0062114} Relasi $\sim$ yang didefinisikan di atas merupakan relasi ekuivalensi.
\end{prop}

\begin{notn} Untuk relasi ekuivalensi $\sim$ yang didefinisikan dalam~\ref{0062111}, kelas ekuivalensi
yang memuat pasangan $(a,b)$ akan dinotasikan dengan % SOURCE-CORRECTION: CH17-C009
 \index{<bracpnt@$\langle a,b \rangle$ (pasangan dalam konstruksi Grothendieck)}%
$\langle a,b \rangle$, bukan dengan~$[(a,b)]$.
\end{notn}

\begin{defn} Misalkan $(S,+)$ suatu semigrup komutatif. Pada $G(S)$, definisikan operasi biner
 \index{penjumlahan!pada grup Grothendieck}%
(yang juga dinotasikan dengan $+$) sebagai berikut:
  \[ \langle a,b \rangle + \langle c,d \rangle := \langle a + c,b + d\rangle\,. \]
\end{defn}

\begin{prop}\label{0062119} Operasi $+$ yang didefinisikan di atas terdefinisi dengan baik, dan di bawah
operasi ini $G(S)$ menjadi grup Abelian. % SOURCE-CORRECTION: CH17-C010
\end{prop}

Grup Abelian $(G(S),+)$ disebut
 \index{G@$G(S)$ (grup Grothendieck dari $S$)}%
 \index{Grothendieck!grup}%
 \index{grup!Grothendieck}%
\df{grup Grothendieck} dari~$S$.

\begin{prop}\label{0062124} Untuk suatu semigrup $S$ dan sembarang $a \in S$, definisikan
pemetaan
 \index{gamma@$\sbsb{\gamma}S$ (pemetaan Grothendieck)}%
   \[ \sbsb{\gamma}S \colon S \sto G(S) \colon s \mapsto \langle s + a, a \rangle\,. \]
Pemetaan $\sbsb{\gamma}S$, yang disebut
 \index{Grothendieck!pemetaan}%
\df{pemetaan Grothendieck}, terdefinisi dengan baik dan merupakan homomorfisme semigrup.
\end{prop}

Pernyataan bahwa pemetaan Grothendieck terdefinisi dengan baik berarti bahwa definisinya tidak bergantung
pada pilihan~$a$. Kita sering hanya menulis $\gamma$ untuk~$\sbsb{\gamma}S$.

\begin{exam}\label{0062126} Baik $\N = \{1,2,3,\dots\}$ maupun $\Z^+ = \{0,1,2,\dots\}$ merupakan
semigrup komutatif di bawah penjumlahan. Keduanya menghasilkan grup Grothendieck yang sama,
   \[G(\N) = G(\Z^+) = \Z\,. \]
\end{exam}

Tidak ada yang menjamin bahwa grup Grothendieck dari suatu semigrup sembarang akan sangat menarik.

\begin{exam}\label{0062127} Misalkan $S$ semigrup aditif komutatif $\Z^+ \cup \{\infty\}$. Maka $G(S) = \{\vc 0\}$.
\end{exam}

\begin{exam}\label{0062128} Misalkan $\Z_0$ semigrup multiplikatif (komutatif) bilangan bulat tak nol.
Maka $G(\Z_0) = \Q_0$, yaitu grup multiplikatif Abelian dari bilangan rasional tak
nol.
\end{exam}

\begin{prop}\label{0062141} Jika $S$ suatu semigrup komutatif, maka
   \[ G(S) = \{\gamma(s) - \gamma(t)\colon s,t \in S\}\,. \]
\end{prop}

\begin{prop}\label{0062142} Jika $r$, $s \in S$, dengan $S$ suatu semigrup komutatif, maka $\gamma(r) = \gamma(s)$ jika dan hanya jika
terdapat $t \in S$ sedemikian sehingga $r + t = s + t$.
\end{prop}

\begin{defn} Suatu semigrup komutatif $S$ memiliki
 \index{sifat pembatalan}%
\df{sifat pembatalan} jika setiap kali $r$, $s$, $t \in S$ memenuhi $r + t = s + t$, berlaku $r = s$.
\end{defn}

\begin{cor}\label{0062144} Misalkan $S$ suatu semigrup komutatif. Pemetaan Grothendieck $\sbsb\gamma S\colon S \sto G(S)$ injektif
jika dan hanya jika $S$ memiliki sifat pembatalan.
\end{cor}

Proposisi berikut menyatakan sifat universal grup Grothendieck.

\begin{prop}\label{0062151} Misalkan $S$ suatu semigrup (aditif) komutatif dan $G(S)$ grup Grothendieck-nya.
 \index{universal!sifat!pemetaan Grothendieck}%
 \index{Grothendieck!pemetaan!sifat universal}%
Jika $H$ suatu grup Abelian dan $\phi\colon S \sto H$ suatu pemetaan aditif, maka terdapat
homomorfisme grup tunggal $\psi\colon G(S) \sto H$ sedemikian sehingga diagram berikut
komutatif.
 \begin{equation}\label{0062151i}
   \xy
     \qtriangle/{>}`{>}`{-->}/[S`\abs{G(S)}`\abs{H};\gamma`\phi`\abs{\psi}]
     \morphism(1100,500)|r|/{-->}/<0,-500>[G(S)`H;\psi]
   \endxy
 \end{equation}
\end{prop}

Dalam diagram sebelumnya, $\abs{G(S)}$ dan $\abs{H}$ hanyalah $G(S)$ dan $H$ yang dipandang sebagai
semigrup, sedangkan $\abs{\psi}$ adalah homomorfisme semigrup yang bersesuaian. Dengan kata lain,
funktor pelupa $\abs{\hphantom{G}}$ ``melupakan'' hanya identitas dan invers, tetapi tidak
operasi penjumlahan. Jadi, segitiga di sebelah kiri merupakan diagram komutatif dalam
kategori semigrup dan homomorfisme semigrup.

\begin{prop}\label{0062154} Misalkan $\phi\colon S \sto T$ homomorfisme semigrup
komutatif. Maka pemetaan $\sbsb{\gamma}T \circ \phi\colon S \sto G(T)$ aditif. Berdasarkan
proposisi~\ref{0062151}, terdapat homomorfisme grup tunggal $G(\phi)\colon G(S)
\sto G(T)$ sedemikian sehingga diagram berikut komutatif.
   \[ \xy
          \Square[S`T`G(S)`G(T);\phi`\sbsb{\gamma}S`\sbsb{\gamma}T`G(\phi)]
      \endxy \]
\end{prop}

\begin{prop}\label{0062155} Pasangan pemetaan $S \mapsto G(S)$, yang memetakan semigrup komutatif ke
 \index{funktor!Grothendieck}%
 \index{Grothendieck!konstruksi!sifat funktorial}%
grup Grothendieck yang bersesuaian, dan $\phi \mapsto G(\phi)$, yang memetakan homomorfisme semigrup ke
homomorfisme-homomorfisme grup (sebagaimana didefinisikan dalam~\ref{0062154}), merupakan funktor kovarian dari kategori % SOURCE-CORRECTION: CH17-C011
semigrup komutatif dan homomorfisme semigrup ke kategori grup Abelian
dan homomorfisme grup.
\end{prop}

Salah satu keuntungan kecil dari penyertaan funktor pelupa yang agak bertele-tele dalam
diagram~\eqref{0062151i} adalah bahwa hal itu memungkinkan kita memandang pemetaan Grothendieck
$\gamma\colon S \mapsto \sbsb{\gamma}S$ sebagai transformasi alami antarfungtor.

\begin{cor}[Kealamian Pemetaan Grothendieck]\label{0062157} Misalkan $\abs{\hphantom{G}}$ funktor pelupa
pada grup Abelian yang ``melupakan'' identitas dan invers, tetapi tidak
operasi grup, seperti dalam~\ref{0062151}. Maka $\abs{G(\hphantom{S})}$ merupakan funktor kovarian dari
kategori semigrup komutatif dan homomorfisme semigrup ke dirinya sendiri.
 \index{Grothendieck!pemetaan!kealamian}%
 \index{alami!transformasi!pemetaan Grothendieck sebagai suatu}%
Selanjutnya, pemetaan Grothendieck $\gamma\colon S \mapsto \sbsb{\gamma}S$ merupakan transformasi
alami dari funktor identitas ke funktor~$\abs{G(\hphantom{S})}$.
  \[ \xy
          \Square[S`T`\abs{G(S)}`\abs{G(T)};\phi`\sbsb{\gamma}S`\sbsb{\gamma}T`\abs{G(\phi)}]
      \endxy \]
\end{cor}
